	The abstract is a concise summary of the dissertation, encapsulating its core content in a clear, logically coherent manner to engage readers. It should be limited to one page, approximately 300–500 words, and avoid including figures or tables.
    
\textbf{Content Elements:}
\begin{itemize}
    \item Research Motivation and Objectives: Describe the background, core research question, and primary goals.
    \item Research Scope: Clearly define the topics and boundaries of the study.
    \item Methodology: Summarize the methods, techniques, or tools used and briefly justify their selection.
    \item Key Results: Highlight the most significant findings or data.
    \item Comparison and Significance: Compare results with the initial objectives and existing literature, emphasizing contributions.
    \item Impact and Future Directions: Discuss the study’s benefits to academic or practical fields and suggest future research or applications.
\end{itemize}


\textbf{Writing Tips:}
\begin{itemize}
    \item Use precise, accessible language to ensure clarity for non-experts.
    \item Draft the abstract after completing other sections to ensure it fully reflects the dissertation’s content.
    \item Format in plain text (no HTML or special formatting) to comply with library archiving and online search requirements.
\end{itemize}

\medskip
\textbf{Keywords:} database, HTML, library archiving, dissertation, methodology